\section{Часть 1. Оценка защищённости разработанного приложения}
\label{sec:part1}

\textbf{Краткое описание приложения.}
Разработано веб-приложение на языке Python с каркасом Flask и встроенной СУБД SQLite. Оно позволяет зарегистрировать пользователя, войти по паролю, выполнять создание, изменение, удаление и поиск двух классов записей и выйти из сессии. Неконфиденциальные заметки хранятся открытым текстом. Конфиденциальные секреты после ввода шифруются схемой Fernet: в базе остаётся шифртекст, на списке отображается только заголовок. Ключ шифрования и секрет приложения записываются в каталог экземпляра с правами только для владельца. Сервер по умолчанию принимает соединения на локальном интерфейсе. Поток данных показан на рисунке~\ref{fig:flow}. Интерфейс~\mdash  на рисунках~\ref{fig:ui} и~\ref{fig:secrets}.

\begin{figure}[H]
	\centering
	\includegraphics[width=0.92\textwidth]{01_flow.png}
	\caption{Разделение заметок и секретов при записи в хранилище}
	\label{fig:flow}
\end{figure}

\begin{figure}[H]
	\centering
	\includegraphics[width=0.88\textwidth]{02_ui.png}
	\caption{Публичная страница и вход в систему}
	\label{fig:ui}
\end{figure}

\begin{figure}[H]
	\centering
	\includegraphics[width=0.88\textwidth]{03_secrets.png}
	\caption{Список секретов без открытого текста тела}
	\label{fig:secrets}
\end{figure}

Запуск сервера: виртуальное окружение Python, установка зависимостей из файла требований, затем \texttt{python app.py}. Служба отвечает на \texttt{http://127.0.0.1:5050/}.

\textbf{Анализ защищённости приложения.}
С позиции операционной системы процесс работает без прав суперпользователя. Каталог экземпляра имеет режим 700, файлы ключа и базы~\mdash  600. Это закрывает доступ другим учётным записям ОС, но не защищает от любого процесса того же пользователя: ключ лежит рядом с базой, и их совместное чтение даёт все секреты. Мера: отдельный системный пользователь и шифрование раздела.

С позиции компонентов и взаимодействия пароли хранятся как хэш PBKDF2-SHA256; изменяющие запросы проверяют CSRF-токен сессии (листинг~\ref{lst:csrf}); выборка записей фильтруется по идентификатору владельца (листинг~\ref{lst:owner}); доступ к СУБД идёт через ORM; шаблоны экранируют вывод. На вход поставлено ограничение частоты неудачных попыток. Заголовки безопасности запрещают встраивание страницы в кадр и подключают CSP. Ограничения: один ключ Fernet на всё приложение; заголовки секретов остаются открытым текстом.

Сетевая политика: прослушивание только \texttt{127.0.0.1}; cookie с флагами HttpOnly и SameSite=Lax; флаг Secure включается при работе через HTTPS. Без TLS вынос порта в сеть позволяет перехватить сессию.

\textbf{Анализ защищённости конфиденциальной информации.}
К конфиденциальным отнесены пароли, тела секретов, ключ Fernet и идентификатор сессии; к неконфиденциальным~\mdash  заголовки и тексты заметок владельца. Разделение сделано разными сущностями и разным способом хранения. Расшифрование тела выполняется только для владельца на экране правки. Поиск секретов идёт по заголовку и не вызывает расшифрование. Ключ нельзя помещать в систему контроля версий.

\textbf{Демонстрация защищённости.}
Практическая проверка: запрос к чужому идентификатору секрета возвращает отказ; POST без токена CSRF отклоняется; в SQLite заметка <<Повестка>> читается как текст, тело секрета~\mdash  двоичный шифртекст, пароль~\mdash  строка PBKDF2. Вид хранилища~\mdash  таблица~\ref{tab:store}.

\begin{table}[ht]
	\begin{tabularx}{\textwidth}{|l|X|}
		\hline
		\textbf{Поле} & \textbf{Значение} \\
		\hline
		заметка & заголовок <<Повестка>>, тело <<встреча в 10:00>> \\
		\hline
		секрет PIN & длина BLOB 100~байт, префикс hex \texttt{674141414141427170485246} \\
		\hline
		alice & \texttt{pbkdf2:sha256:1000000\$6umpCH}\ldots \\
		\hline
		bob & \texttt{pbkdf2:sha256:1000000\$upbYko}\ldots \\
		\hline
		права & каталог экземпляра 0700; \texttt{fernet.key}, \texttt{notes.db}, \texttt{flask.secret} 0600 \\
		\hline
	\end{tabularx}
	\caption{Вид записей в SQLite и режимы файлов экземпляра}
	\label{tab:store}
\end{table}

\begin{lstlisting}[caption={Проверка CSRF-токена изменяющих запросов},label={lst:csrf}]
@app.before_request
def _csrf_protect():
    if request.method in ("POST", "PUT", "DELETE", "PATCH"):
        token = session.get("_csrf")
        sent = request.form.get("_csrf") or \
            request.headers.get("X-CSRF-Token")
        if not token or not sent or not secrets.compare_digest(
            str(token), str(sent)
        ):
            abort(400, description="CSRF token missing or invalid")
\end{lstlisting}

\begin{lstlisting}[caption={Выборка секрета только по идентификатору владельца},label={lst:owner}]
secret = Secret.query.filter_by(
    id=secret_id, user_id=current_user().id
).first_or_404()
\end{lstlisting}
